\chapter{\v{S}pecifik\'{a}cia a n\'{a}vrh syst\'{e}mu}
\label{ch:spec}

V tejto kapitole vyslovíme požiadavky na systém a navrhneme riešenie.

\section{\'{U}vod do problematiky}
\label{sec:intro}

Majme projekt skladajúci sa z väčšieho množstva zdrojových súborov, ktoré treba
skompilovať (alebo pre webové aplikácie minifikovať) predtým, ako z nich vytvoríme
výslednú aplikáciu. Keďže tento postup je presne daný a častokrát opakovaný,
je veľmi výhodné ušetriť čas potrebný na písanie všetkých príkazov za sebou
použitím nejakého systému, ktorý to zautomatizuje. Okrem automatizácie chceme
určite ušetriť aj čas pri takzvanej rekompilácii, aby sa po zmene niekoľkých súborov
nemusela znova kompilovať celá aplikácia. Niekoľko existujúcich systémov spomíname
v časti~\ref{sec:existing}.

V prípade spoločnosti pracujúcej na projektoch obsahujúcich rádovo stovky až tisíce
súborov je takýto systém nevyhnutný. Avšak aj pri použití automatizácie a použitia
výsledkov predchádzajúcich behov môže rekompilácia, najmä pri väčších zmenách,
trvať dlho, čo znižuje nielen produktivitu zamestnancov, ale aj zisky spoločnosti.
Preto niektoré korporácie využívajú na kompiláciu (ktorá je pri veľkom počte
nezávislých súborov inherentne vysoko paralelizovateľná) nielen pracovné stanice
zamestnancov, ale, v prípade internetových gigantov aj ich serverové farmy.

Avšak narážame na problém. V čase písania tejto práce sú všetky takto distribuované
riešenia uzavreté, neprístupné verejnosti, iba spoločnosti, ktorá ho vyvinula
a drží si ho pod svojou strechou. Toto je veľká nevýhoda pre začínajúce start-up
organizácie, ktoré, aj keď majú k dispozícii servery, ktoré by takto mohli využiť,
nemôžu si dovoliť vývoj takéhoto riešenia pre vlastnú potrebu.

\section{Formul\'{a}cia probl\'{e}mu}
\label{sec:problem}

Sformulujme teda problém na základe predchádzajúcej časti a pridajme ďalšie
požiadavky:

\begin{itemize}
  \item Automatizácia kompilácie a linkovania zdrojových súborov do výslednej aplikácie.
  \item Využívanie predchádzajúcich výsledkov pri ďalších behoch.
  \item Distribúcia kompilácie na servery, pokiaľ sú dostupné.
  \item Sprístupnenie existujúcich behov a výsledkov celej organizácii.
  \item Rozšíriteľná podpora rôznych jazykov a postupov.
  \item Minimalizovanie potreby ručnej konfigurácie.
  \item Podpora VCS\footnote{Version Control System}, najmä Git.\footnote{\url{https://git-scm.com}}
  \item Vydané pod open-source licenciou.
\end{itemize}

Pozrime sa podrobnejšie na niektoré požiadavky.

\subsubsection{Sprístupnenie existujúcich behov a výsledkov celej organizácii}
Ak dvaja zamestnanci pracujú na projekte potrebujúcom kompiláciu tisícky súborov,
nie je dôvod aby obaja, pri prvom pokuse skompilovať takýto projekt, museli čakať
niekoľko hodín pred tým ako môžu ďalej pracovať. Preto plánujeme vytvoriť distribuované
úložisko objektov vytvorených kompiláciou na serveroch, ktoré budú takto dostupné
všetkým používateľom daného systému. Toto všetko za účelom čo najväčšieho ušetrenia
výpočtového výkonu a času spotrebovaného kompiláciou.

\subsubsection{Rozšíriteľná podpora rôznych jazykov a postupov}
V prvej verzii tohto systému plánujeme implementovať podporu pre niekoľko najpoužívanejších
jazykov. Avšak uvedomujeme si, že takýto systém by mal byť nezávislý na použitých
nástrojoch a preto jeho súčasťou bude aj jednoduchý spôsob pridávania podpory
pre iné platformy.

\subsubsection{Minimalizovanie potreby ručnej konfigurácie}
V prípade existujúcich riešení je častokrát až zarážajúce množstvo konfigurácie
ktorú treba vykonať pri použití takého systému v projekte. Náš systém by mal
tieto záležitosti čo najviac zjednodušiť.

\subsubsection{Podpora VCS}
Veľké množstvo projektov závisí na súboroch, knižniciach, ktoré sa nenachádzajú
v ich repozitári, ale sú dostupné pod nejakým VCS\@. Náš systém by mal, pri
vhodnej špecifikácii dostupnosti týchto závislostí v konfigurácii, byť schopný
automaticky stiahnuť a použiť pri kompilácii takúto závislosť, bez potreby akejkoľvek
činnosti od používateľa.

\section{Existuj\'{u}ce syst\'{e}my}
\label{sec:existing}
V tejto časti popíšeme niekoľko open-source build systémov. Vo všetkých prípadoch
ide ale iba o lokálnu kompiláciu.

\subsubsection{GNU Make}
Najstarší a stále pravdepodobne najpoužívanejší build systém. Využíva súbor
nazvaný Makefile, ktorý obsahuje popis pravidiel na vytváranie súborov na základe
ich prípony. Na zabránenie opakovanej kompilácii porovnáva časy posledného zápisu
do súborov (zdrojového a výsledného), čo ale nezabraňuje opakovanej kompilácii
napríklad v prípade obnovenia zdrojového súboru zo zálohy. Na zjednodušenie písania
pravidiel v Makefile vzniklo niekoľko systémov, ktoré tvorbu Makefile automatizujú.
Medzi najpoužívanejšie patria GNU Autotools a CMake.

\subsubsection{Apache Ant}
Apache Ant je build systém, ktorý adresuje niekoľko problémov s prenositeľnosťou
Makefile súborov medzi platformami, keďže tieto špecifikujú akcie ako postupnosť
príkazov v shelli. Konfigurácia je vykonávaná v XML súboroch, ktoré definujú
strom úloh (tasks) pre projekt. Každá z úloh je implementovaná ako trieda v
jazyku Java, v ktorom je implementovaný aj celý Ant a teda je jednoduché
pridať nové úlohy, ako aj preniesť existujúce úlohy na nové systémy. Tento systém
využívajú najmä projekty písane v jazyku Java.

\subsubsection{Bazel}
Open-source verzia interného build systému spoločnosti Google, ktorý bol prvý krát
zverejnený v roku 2014. Využíva podobné koncepty ako Apache Ant. Funkcionalita tohto
systému bola popísaná v roku 2012 sériou článkov na Google Developers blogu
\cite{BlazeDesign} a slúži ako inšpirácia pri návrhu nášho systému. Podľa plánov
vydania \cite{BazelRoadmap} sa ale Google nechystá zverejniť ich implementáciu
distribuovanej kompilácie, ktorá je popísaná vo vyššie spomínanej sérii článkov.

\section{N\'{a}vrh rie\v{s}enia}
\label{sec:solution}
V tejto časti popíšeme jednotlivé komponenty systému a ich funkcionalitu.

\subsection{Master}
\label{sub:solution:master}
Aplikáciu bude predstavovať program, bežiaci na užívateľskom počítači. Tento program
nazvime {\it master}.

Úlohou mastera bude v okamihu zavolania a špecifikovania cieľa kompilácie spracovať
konfiguráciu cieľov daného projektu a vytvoriť strom závislostí. Následne master
kontaktuje jednotlivé {\it slave} komponenty bežiace na rôznych strojoch (lokálny
počítač, serverová farma a pod.) a rozdistribuuje ciele na tieto počítače čo možno
najefektívnejšie.

V prípade, že cieľ alebo podciele závisia na vonkajšom projekte, tento projekt je,
v prípade, že jeho repozitár je dostupný, stiahnutý a použitý. V prípade nedostupnosti
niektorej zo závislostí program ukončí beh s chybovou hláškou.

\subsection{Slave}
\label{sub:solution:slave}
Slave je program starajúci sa o vykonávanie úloh zadaných masterom. Na fyzickom počítači,
či už lokálnom, alebo na serveri bude bežať zvyčajne niekoľko slave programov, nakoľko
jeden slave program bude navrhnutý na vykonávanie úloh pre práve jeden cieľ. Keď voľný
slave obdrží úlohu, v prvom kroku sa pozrie do zdieľaného úložiska objektov, ak je
toto dostupné, či už žiadaný súbor neexistuje. V prípade jeho existencie slave
pošle výsledný súbor priamo masterovi a nebude vykonávať kompiláciu. V prípade
neexistencie výsledného súboru v žiadanej verzii bude cieľ preložený na postupnosť
akcií a tieto budú vykonané.

\subsection{\'{U}lo\v{z}isko objektov}
\label{sub:solution:storage}
Úložisko objektov slúži ako cache na zabránenie opakovanej kompilácii rovnakých
súborov. Výsledné súbory tu budú nahrané po dokončení kompilácie slave programom.
Ukladanie bude prebiehať podľa názvu zdrojového súboru a jeho kontrolného súčtu,
aby sa ľahko rozlíšili rôzne verzie súboru, napríklad lokálne modifikovaná od
tzv. HEAD verzie\footnote{najnovšia verzia nahraná zverejnená v hlavnom repozitári}.

V prípade nášho úložiska sme sa rozhodli použiť na kontrolné súčty hashovaciu funkciu
SHA512, kvôli čo najväčšej minimalizácii kolízií.

\subsection{Ciele kompil\'{a}cie}
\label{sub:solution:targets}
Ciele kompilácie delíme nasledovne:

\begin{itemize}
  \item Súbory, delia sa na 2 podskupiny:
    \begin{itemize}
      \item Zdrojové súbory
      \item Generované súbory
    \end{itemize}
  \item Pravidlá
\end{itemize}

Pri vyžiadaní kompilácie cieľa vytvoríme orientovaný graf závislostí, ktorý vyžadujeme,
aby bol acyklický. V prípade nájdenia cyklu systém ukončí beh s chybou o nesplniteľnosti
požiadavky. Takto zostrojený graf bude, na základe vyššie uvedených typov cieľov,
bipartitný.

Nasleduje bližší popis jednotlivých typov cieľov:

\subsubsection{Zdrojové súbory}
Toto sú súbory, ktoré sa nachádzajú priamo v repozitári projektu, sú vytvorené
ľuďmi (alebo nástrojmi) a nie generované týmto systémom. V grafe závislostí
tvoria listy grafu.

\subsubsection{Generované súbory}
Súbory vytvorené ako výstup z nejakého pravidla, častokrát slúžia zároveň aj ako
vstupné súbory pre ďalšie akcie.

\subsubsection{Pravidlá}
Popisujú vzťah a postup vytvorenia výstupného súboru zo súborov vstupných.
Každé pravidlo má práve jeden výstupný súbor a 0 alebo viac vstupných. Pravidlá
zahŕňajú postupy ako kompilácia, linkovanie a podobne.

\begin{figure}[h]
  \centerline{\includegraphics[width=0.6\textwidth]{images/dependency_graph}}
  \caption[Graf závislostí]{Príklad grafu závislostí, okrúhle vrcholy predstavujú
  súbory, štvorcové vrcholy predstavujú pravidlá.}
  \label{img:dependency_graph}
\end{figure}

Takto zostrojený graf bude potom použitý slave komponentmi na vytvorenie výstupných
súborov s použitím pravidiel a príkazov ktoré so sebou nesú.

\newpage

\subsection{Vo\v{l}ba technol\'{o}gi\'{i}}
\label{sub:solution:tech}

Pre tento projekt sme zvolili 2 programovacie jazyky, jazyk Go\footnote{\url{https://golang.org}}
pre používateľskú (master) časť aplikácie a Erlang\footnote{\url{https://www.erlang.org}}
pre serverovú (slave) časť. Objektové úložisko bude využívať program rsync\footnote{\url{https://rsync.samba.org}},
ktorý dokáže pristupovať k súborom pomocou rôznych protokolov. Na strane servera bude úložisko
a jeho distribuovanosť prenechaná na súborový systém, kde existuje niekoľko vhodných
open-source alternatív. Tieto technológie budú bližšie popísané v nasledujúcej kapitole.

Zdrojový kód bude verzionovaný systémom Git\footnote{\url{https://git-scm.com}} a
zverejnený pod licenciou MIT v GitHub repozitári na adrese \url{https://github.com/imterra/forge}.
